% Generated by scripts/build_readers.py; change the reviewed sources and rebuild.
\def\EditionMode{0}
\def\EditionLabel{Diplomatic reader}
\def\EditionStatus{Complete checked edition}
\documentclass[10pt,oneside]{book}
\usepackage{fontspec}
\setmainfont{Latin Modern Roman}
\setsansfont{Latin Modern Sans}
\usepackage{amsmath,amssymb,mathtools}
\usepackage{unicode-math}
\setmathfont{Latin Modern Math}
\usepackage{xcolor,tikz-cd,lua-ul,marginnote,pdflscape}
\usetikzlibrary{arrows.meta,decorations.pathmorphing,calc,shapes.geometric}
\tikzcdset{ampersand replacement=\&}
\tikzcdset{every label/.append style={text=black}}
\usepackage[a4paper,left=18mm,right=45mm,top=18mm,bottom=22mm,marginparwidth=28mm,marginparsep=7mm]{geometry}
\usepackage[unicode,colorlinks=true,urlcolor=black,linkcolor=black]{hyperref}
\usepackage{bookmark}
\setlength{\parindent}{0pt}
\setlength{\parskip}{4pt plus 1pt minus 1pt}
\setlength{\emergencystretch}{1.5em}
\raggedbottom
% Include final-line descenders inside the declared text-height rectangle.
\maxdepth=0pt
\definecolor{PSHand}{RGB}{25,58,132}
\definecolor{PSDelete}{RGB}{146,37,37}
\definecolor{PSTypedOld}{RGB}{95,95,95}
\ifnum\EditionMode=2\colorlet{blue}{black}\else\colorlet{blue}{PSHand}\fi
\RenewDocumentCommand{\textcircled}{m}{\tikz[baseline=(PSCircle.base)]\node[draw,circle,inner sep=.5pt,line width=.4pt](PSCircle){#1};}
\newif\ifPSInApparatus
\newcounter{psnote}
\newcounter{psvariant}
\newcounter{pslinkmark}
\newif\ifPSInLink
\newsavebox{\PSNoteBox}
\newdimen\PSNoteLimit
{\footnotesize\global\PSNoteLimit=3\baselineskip}
\newcommand{\CurrentScanURL}{}
\newcommand{\CurrentKey}{}
\newcommand{\CurrentScan}{}
\newcommand{\pssource}[3]{\psflushmargins\clearpage\gdef\CurrentKey{#1}\gdef\CurrentScan{#2}\gdef\CurrentScanURL{#3}\setcounter{psvariant}{0}\typeout{PS-SOURCE: #1}}
\newcommand{\psdivision}[2]{\pdfbookmark[0]{#1}{#2}\par\noindent{\large\bfseries #1\par}\medskip}
\newcommand{\pslinkmarker}[2]{\stepcounter{pslinkmark}\latelua{texio.write_nl('log','PS-LINK-SHIP|#1|\CurrentKey|\arabic{pslinkmark}|#2')}}
\newcommand{\pslink}[2]{\ifPSInLink\pslinkmarker{#1}{nested}#2\else\begingroup\PSInLinktrue\href{\CurrentScanURL}{\pslinkmarker{#1}{outer}#2}\endgroup\fi}
\newcommand{\scanpage}[1]{\ifnum\EditionMode=2\leavevmode\marginpar{\scriptsize\sffamily\pslink{scanpage}{scan #1}}\else\par\noindent{\small\sffamily\pslink{scanpage}{[scan #1]}}\par\medskip\fi}
\newcommand{\mspage}[1]{\ifnum\EditionMode=2\leavevmode\marginpar{\scriptsize\sffamily p.\,#1}\else{\small\sffamily[p.\,#1]}\fi}
\newcommand{\hand}[1]{{\ifnum\EditionMode=2\color{black}\else\color{PSHand}\fi #1}}
\let\over\relax
\NewDocumentCommand{\over}{O{} m m}{{\textcolor{PSTypedOld}{#2}\ensuremath{^{\!\nearrow}}\if\relax\detokenize{#1}\relax\hand{#3}\else #3\fi}}
\NewDocumentCommand{\inserted}{O{} m}{{\if\relax\detokenize{#1}\relax\hand{\ensuremath{\langle}#2\ensuremath{\rangle}}\else\ensuremath{\langle}#2\ensuremath{\rangle}\fi}}
\newcommand{\deleted}[1]{{\color{PSDelete}\ifmmode\text{[del.: }#1\text{]}\else[\textsf{del.:} #1]\fi}}
\newcommand{\xout}[1]{{\color{PSTypedOld}\ifmmode\text{[x: }#1\text{]}\else[\textsf{x:} #1]\fi}}
\let\PSMathUnderline\underline
\RenewDocumentCommand{\underline}{+m}{\ifmmode\PSMathUnderline{#1}\else\underLine{#1}\fi}
\NewDocumentCommand{\uline}{+m}{\underline{#1}}
\NewDocumentCommand{\uuline}{+m}{\ifmmode\PSMathUnderline{\PSMathUnderline{#1}}\else\underLine{\underLine{#1}}\fi}
\newunderlinetype\PSWaveLine{\leaders\hbox{\lower.7ex\hbox{\tikz[x=1pt,y=1pt,baseline=0pt]{\draw[line width=.3pt] (0,0) .. controls (1,1) and (2,1) .. (3,0) .. controls (4,-1) and (5,-1) .. (6,0);}}}}
\NewDocumentCommand{\uwave}{+m}{{\PSWaveLine #1}}
\newcommand{\unsure}[1]{\pslink{unsure}{[#1\ifmmode\,?\else\,?\fi]}}
\newcommand{\illegible}[1]{\ifmmode\text{\pslink{illegible}{[illeg.\ #1]}}\else\pslink{illegible}{[illeg.\ #1]}\fi}
\NewDocumentCommand{\diagram}{+m}{\ifmmode\vcenter{\hbox{#1}}\else\par\medskip{\centering #1\par}\medskip\psflushmargins\fi}
\newenvironment{diagramblock}{\par\begin{center}}{\end{center}\par}
\newcommand{\PSMathNotes}{}
\makeatletter
\NewDocumentCommand{\psplacednote}{m +m}{%
      \sbox{\PSNoteBox}{\parbox{\marginparwidth}{\footnotesize\raggedright\hand{\textsf{m#1.} #2}}}%
      \ifdim\dimexpr\ht\PSNoteBox+\dp\PSNoteBox\relax<\PSNoteLimit
        \leavevmode\marginpar{\usebox{\PSNoteBox}}%
      \else
        \par\nobreak\smallskip\hrule\nobreak\smallskip
        {\footnotesize\textsf{Marginal note m#1.}\quad\hand{#2}\par}
        \nobreak\smallskip\hrule\smallskip
      \fi}
\newcommand{\psflushmargins}{\ifmmode\else\ifinner\else\ifx\PSMathNotes\@empty\else\let\PSTempNotes\PSMathNotes\global\let\PSMathNotes\@empty\PSTempNotes\fi\fi\fi}
\NewDocumentCommand{\margin}{+m}{%
  \ifPSInApparatus\ifmmode\hand{\text{[marg.: }#1\text{]}}\else\hand{\textsf{[marg.: }#1\textsf{]}}\fi
  \else\stepcounter{psnote}%
    \ifmmode
      \text{\textsuperscript{m\arabic{psnote}}}%
      \begingroup\edef\PSQueue{\endgroup\noexpand\g@addto@macro\noexpand\PSMathNotes{\noexpand\psplacednote{\arabic{psnote}}{\unexpanded{$#1$}}}}\PSQueue
    \else
      \textsuperscript{m\arabic{psnote}}\psplacednote{\arabic{psnote}}{#1}%
    \fi
  \fi}
\AtBeginDocument{\everymath\expandafter{\the\everymath\aftergroup\psflushmargins}\everydisplay\expandafter{\the\everydisplay\aftergroup\psflushmargins}}
\makeatother
\newcommand{\psnumber}[2]{\expandafter\gdef\csname psnum@#1\endcsname{#2}}
\NewDocumentCommand{\var}{+m +m m m}{\ifnum\EditionMode=0 #1\else\ifnum\EditionMode=1\ifmmode\text{\textsuperscript{\tiny\csname psnum@#4\endcsname}}\else\textsuperscript{\tiny\csname psnum@#4\endcsname}\fi\fi #2\fi}
\newcommand{\psapparatusstart}{\par\bigskip\hrule\smallskip\PSInApparatustrue\begingroup\footnotesize\raggedright\textsf{Apparatus for scan \CurrentScan}\par}
\newcommand{\psapparatusend}{\par\endgroup\PSInApparatusfalse}
\NewDocumentCommand{\psentry}{m m m +m +m m}{\par\noindent\textbf{#2. #3.}\quad\textit{Typescript:} #4\quad\textit{Corrected:} \if\relax\detokenize{#5}\relax\textnormal{[omitted]}\else #5\fi\quad #6\par}
\newcommand{\psfrontmatter}{%
\hypersetup{pageanchor=false}\begin{titlepage}\raggedright
{\Large Alexander Grothendieck\par}\vspace{12mm}
{\Huge Pursuing Stacks\par}\medskip{\LARGE 1983\par}\vspace{12mm}
{\Large\EditionLabel\par}\vspace{8mm}
{\large\EditionStatus\par}\vfill
Université de Montpellier, Grothendieck archive, cotes 134-2--134-8.\par
Including the 26 source pages of R. Brown's translation at cote 134-2, sites 115--143.\par
\vspace*{2mm}\end{titlepage}\hypersetup{pageanchor=true}
\chapter*{Editorial statement}
AI-transcribed by Claude Opus 5.5 (Anthropic, effort high); checked against the scans, and this edition built, by GPT-6 Astra (OpenAI, ChatGPT Pro mode); prompts and records by Claude (Anthropic).

The witness consulted and not reproduced is the LaTeX edition of Mateo Carmona with the collaboration of Ulrik Buchholtz (arXiv 2111.01000 v2, 2021), extended from the 2015 edition by ``the scrivener'', CC0 1.0.

The translation in the appendix is R. Brown's, prepared with assistance from L. Breen and J-L. Loday, as his note says; the rights in it belong to its translator or his heirs.

The rights in the original work belong to its author or the author's representatives or heirs. Where we hold rights, in the AI-produced transcription, translation, corrections and typesetting, they are dedicated to the public domain under CC0 1.0. This edition is published for educational, non-profit purposes, and we claim no rights we do not hold. If you hold rights in this work and object to its publication here, please open a GitHub issue and we will take it down.

\section*{How to read this form}
\ifnum\EditionMode=0
The diplomatic form preserves typed and handwritten wording, cancellations, insertions, overwritten readings, underlining, doubts, and original page numbers. Blue marks handwriting; red labels handwritten cancellations; gray labels typewriter cancellations and underlying overwritten text. Angle brackets mark insertions; a superscript arrow separates an overwritten reading from its replacement. Typed insertions and typed overwriting remain plain. The disposition of long marginal notes is identified explicitly.
\else
The corrected text applies the author's insertions, replacements, cancellations, and recorded transpositions. It retains his English and notation, including consistent forms such as ``standart'', ``developped'', ``it's'' for ``its'', and ``A la''. An isolated slip is corrected only when specifically recorded and justified. Mathematical anomalies are preserved and discussed in the guide.
\ifnum\EditionMode=1
Blue continues to identify handwriting. A numbered call identifies each maximal editorial operation, with its full diplomatic and corrected readings and reason in the apparatus following the source page. Internal layers of a wholly replaced or deleted passage are preserved within that entry rather than duplicated as separate entries. These entries are generated from the same variants that govern the running text.
\else
This reading form omits the apparatus, cancelled material, superseded readings, insertion brackets, and color distinctions. It retains the resulting text, marginal notes, underlining, and unresolved doubts. Scan and source-page references appear in the margin.
\fi\fi

Every scan mark and remaining doubt links to its source page in the archive viewer. Doubts are editorial reading marks, not emendations. The source-page breaks identify the archival sequence; a source page may occupy more than one printed page. A partial working edition is explicitly labelled and contains only reviewed source pages.
\clearpage}

\begin{document}
\pssource{ps_134-4_042}{134-4/42}{https://grothendieck.umontpellier.fr/web/viewer.html?file=../134-4.pdf\#page=43}
\psnumber{ps_134-4_042-v0001}{1}
\psnumber{ps_134-4_042-v0002}{2}
\psnumber{ps_134-4_042-v0003}{3}
\psnumber{ps_134-4_042-v0004}{4}
\psnumber{ps_134-4_042-v0005}{5}
\psnumber{ps_134-4_042-v0006}{6}
\psnumber{ps_134-4_042-v0007}{7}
\psnumber{ps_134-4_042-v0008}{8}
\psnumber{ps_134-4_042-v0009}{9}
\psnumber{ps_134-4_042-v0010}{10}
\psnumber{ps_134-4_042-v0011}{11}
\psnumber{ps_134-4_042-v0012}{12}
\psnumber{ps_134-4_042-v0013}{13}
\psnumber{ps_134-4_042-v0014}{14}
\psnumber{ps_134-4_042-v0015}{15}
%% BEGIN SOURCE ps_134-4_042
% Pursuing Stacks (cotes 134-2 to 134-8): archive cote 134-4, site page 42 (PDF page 43), typescript page 130.
% typescript. Draft compared: ps3.tex:2109-2183.
% AI-transcribed by Claude Opus 5.5 (claude-opus-5-5, adaptive thinking, effort high), Message Batches API, batch ps2 (msgbatch_01NVwr1Yuo8qJf5rykeGBXPK), 2026-11-12; image review recorded in records/CHECK_DECISIONS_134_3_4.tsv.
\scanpage{134-4/42}
\mspage{130}

\hand{53.}\quad Before pursuing yesterday's reflection about the $\hand{\pi}_o$-functor and
it's relation to homotopy structures on M, \ it seems more convenient to
interpolate \var{\xout{a reflection}}{}{typewriter deletion}{ps_134-4_042-v0001} some more or less obvious ``functorialities'' on
the homotopy notions just developped. \ They all seem to turn around the
relationship of such notions in M with a more or less arbitrary functor
\[ u : M \hand{\longrightarrow} N , \]
for the time being I am not making any special assumption on \hand{$u$}. \ In terms
of the four ways we got for describing a homotopy structure in M, \ we get
four corresponding natural conditions of ``compatibility'' of \hand{$u$} with a
given homotopy structure h, \ namely :

(i) \ If \ $f \underset{h}{\hand{\sim}} g$ \ in M, \var{\inserted[typed]{i.e. $(f;g)\hand{\in} R_{\hand{h}}$}}{i.e. $(f;g)\hand{\in} R_{\hand{h}}$}{typed insertion}{ps_134-4_042-v0002} \ then \ u(f) = u(g) .

(i') \ If $\underline{I}$ =(I,$\hand{\delta}_o$,$\hand{\delta}_1$)$\hand{\in}\hand{\Sigma}_h$ is a weak homotopy interval, \var{\inserted{then}}{\hand{then}}{insertion}{ps_134-4_042-v0003} u($\hand{\delta}_o$) = u($\hand{\delta}_1$).

(ii) \ If f$\hand{\in}$ W$_h$ i.e. f is a homotopism, \ u(f) is an isomorphism.

(ii') If \ X$\hand{\in}$ C$_h$ is a contractible object, \ u(X) is a final object (NB
here we assume that u(e$_M$) \var{\over{is}{is}}{\hand{is}}{writing over}{ps_134-4_042-v0004}\var{\deleted{\illegible{1}}}{}{deletion}{ps_134-4_042-v0005} a final object in N).

\var{\inserted{(NB.}}{\hand{(NB.}}{insertion}{ps_134-4_042-v0006} It is understood implicitly, \ whenever dealing with intervals and with
contractible \var{\xout{elemen}}{}{typewriter deletion}{ps_134-4_042-v0007} objects, \ that M admits finite products.\hand{)}

The conditions (i) and (ii) are clearly equivalent, \ and equivalent
to the \var{requitement}{requirement}{slip corrected}{ps_134-4_042-v0008} that u factors into
\hand{(1)}
\[ M\hand{\longrightarrow} M_h \hand{\longrightarrow} N , \]
where M$\hand{\rightarrow}$M$_h$ is the canonical functor of M into the corresponding \uline{homotopy}
\uline{-types} category. \ We also have the tautological implications \ (i)$\hand{\Rightarrow}$(i')
and (ii)$\hand{\Rightarrow}$(ii'). \ Moreover we have \var{\inserted{(trivially)}}{\hand{(trivially)}}{insertion}{ps_134-4_042-v0009} (ii')$\hand{\Rightarrow}$(i') \ whenever the homotopy
structure on M is a contractibility structure, \ and (i\unsure{\var{\over{'}{'}}{\hand{'}}{writing over}{ps_134-4_042-v0010}})$\hand{\Longrightarrow}$(i) \ whenever
the homotopy structure on M is a weak homotopy interval structure, \ and
\uline{moreover}\textsuperscript{\hand{\unsure{*}}} \ the functor u commutes to \var{\xout{finite}}{}{typewriter deletion}{ps_134-4_042-v0011}\var{\inserted[typed]{finite}}{finite}{typed insertion}{ps_134-4_042-v0012} products. \ All these implica-%
tion are summarized in the diagram
\margin{$*$}
% Hand-drawn arrows are marked blue; typed objects remain plain.
\diagram{\begin{tikzcd}[arrows=Rightarrow]
(i) \ar[color=blue, r,Leftrightarrow] \ar[color=blue, d] \& (ii) \ar[color=blue, d] \\
(i') \ar[color=blue, u, bend left, "\hand{\ensuremath{\mathrm{Hom}_3}}" '{xshift=12pt}, "\hand{\ensuremath{*}}"] \& (ii') \ar[color=blue, l, "\hand{\ensuremath{\mathrm{Hom}_4}}" '] 
\end{tikzcd}\ \hand{\unsure{\ensuremath{*|1}}}\hand{,}}

where the symbol Hom$_{\hand{3}}$ or Hom$_4$ \ indicates that the implication qualified by
it \var{\over[typed]{is}{is}}{is}{typewriter overwriting}{ps_134-4_042-v0013} valid provided we assume that the homotopy structure \var{\over[typed]{o}{i}}{i}{typewriter overwriting}{ps_134-4_042-v0014}\var{\over[typed]{n}{n}}{n}{typewriter overwriting}{ps_134-4_042-v0015} M is in
Hom$_{\hand{3}}$ \ (namely is defined in terms of weak homotopy intervals) \ resp. in
Hom$_4$ (namely is a contractibility structure), \ and where (*) denotes the
extra assumption of commutation of u with finite products.

\margin{42}

%% END SOURCE ps_134-4_042
\ifnum\EditionMode=1
\psapparatusstart
\psentry{ps_134-4_042-v0001}{1}{typewriter deletion}{\xout{a reflection}}{}{The typewriter strikes out this material; it leaves the corrected running text.}
\psentry{ps_134-4_042-v0002}{2}{typed insertion}{\inserted[typed]{i.e. $(f;g)\hand{\in} R_{\hand{h}}$}}{i.e. $(f;g)\hand{\in} R_{\hand{h}}$}{The typed insertion is incorporated at its indicated position; insertion brackets leave the running text.}
\psentry{ps_134-4_042-v0003}{3}{insertion}{\inserted{then}}{\hand{then}}{The handwritten insertion is incorporated at its indicated position; insertion brackets leave the running text.}
\psentry{ps_134-4_042-v0004}{4}{writing over}{\over{is}{is}}{\hand{is}}{The superposed handwritten reading replaces the underlying reading.}
\psentry{ps_134-4_042-v0005}{5}{deletion}{\deleted{\illegible{1}}}{}{The author cancels this material by hand; it leaves the corrected running text.}
\psentry{ps_134-4_042-v0006}{6}{insertion}{\inserted{(NB.}}{\hand{(NB.}}{The handwritten insertion is incorporated at its indicated position; insertion brackets leave the running text.}
\psentry{ps_134-4_042-v0007}{7}{typewriter deletion}{\xout{elemen}}{}{The typewriter strikes out this material; it leaves the corrected running text.}
\psentry{ps_134-4_042-v0008}{8}{slip corrected}{requitement}{requirement}{Isolated t for r, settled by correctly typed requirement in reviewed 134-4/007.}
\psentry{ps_134-4_042-v0009}{9}{insertion}{\inserted{(trivially)}}{\hand{(trivially)}}{The handwritten insertion is incorporated at its indicated position; insertion brackets leave the running text.}
\psentry{ps_134-4_042-v0010}{10}{writing over}{\over{'}{'}}{\hand{'}}{The superposed handwritten reading replaces the underlying reading.}
\psentry{ps_134-4_042-v0011}{11}{typewriter deletion}{\xout{finite}}{}{The typewriter strikes out this material; it leaves the corrected running text.}
\psentry{ps_134-4_042-v0012}{12}{typed insertion}{\inserted[typed]{finite}}{finite}{The typed insertion is incorporated at its indicated position; insertion brackets leave the running text.}
\psentry{ps_134-4_042-v0013}{13}{typewriter overwriting}{\over[typed]{is}{is}}{is}{The superposed typed reading replaces the underlying reading.}
\psentry{ps_134-4_042-v0014}{14}{typewriter overwriting}{\over[typed]{o}{i}}{i}{The superposed typed reading replaces the underlying reading.}
\psentry{ps_134-4_042-v0015}{15}{typewriter overwriting}{\over[typed]{n}{n}}{n}{The superposed typed reading replaces the underlying reading.}
\psapparatusend
\fi

\end{document}
